\begin{enumerate}[label=\textbf{\ansnum.\arabic*}]
  \item \verb|sudo usermod -aG dialout $USER|. Group memberships are read when the login session starts; a running session keeps the old list, so you must log out and in again (or reconnect SSH).
  \item The kernel numbers USB serial devices in the order in which it detects them (depends on plug order, hub, timing). udev rules give a stable name with \texttt{SYMLINK+=}.
  \item File \texttt{/etc/udev/rules.d/99-lidar.rules}:
\begin{verbatim}
SUBSYSTEM=="tty", ATTRS{idVendor}=="10c4", ATTRS{idProduct}=="ea60", \
  ATTRS{serial}=="0001", SYMLINK+="lidar", GROUP="dialout", MODE="0660"
\end{verbatim}
  Then \verb|sudo udevadm control --reload-rules| and \verb|sudo udevadm trigger|.
  \item (1) \verb|SUBSYSTEM="tty"| uses a single \texttt{=}: this assigns instead of matching, it must be \verb|==|. (2) Only the vendor ID is matched, so \emph{every} tty device of this vendor matches and would get the same link name \texttt{/dev/imu}; add the product ID and, for identical adapters, the serial number.
  \item \verb|lsusb| (shows \texttt{ID vendor:product}), \verb|udevadm info -a -n /dev/DEVICE| (shows \texttt{ATTRS\{idVendor\}} and \texttt{ATTRS\{idProduct\}}), and \verb|sudo dmesg| after plugging in (shows \texttt{idVendor=...}, \texttt{idProduct=...}).
  \item \texttt{ATTR} reads an attribute of the device itself; \texttt{ATTRS} searches the device \emph{and all its parents}. The tty device is a child of the USB device, and the vendor, product and serial attributes belong to the USB parent, so they are found with \texttt{ATTRS}.
  \item \verb|pactl list short sources|, then \verb|pactl set-default-source NAME2|, \verb|pactl set-source-volume @DEFAULT_SOURCE@ 80%|, \verb|pactl set-source-mute @DEFAULT_SOURCE@ 0|.
  \item \verb|arecord -f S16_LE -r 16000 -c 1 -d 5 test.wav| and \verb|aplay test.wav|. Data: $16\,000\cdot2\cdot1\cdot5=160\,000$ bytes.
  \item The sound server runs inside the session of a user; root or a service has none (or a different one). Run the service as the user (\texttt{User=}), set \texttt{XDG\_RUNTIME\_DIR=/run/user/UID}, and enable lingering so the user session exists without a login.
  \item udev kills long-running programs started with \texttt{RUN}, and they run in a restricted environment, before the system is completely up. Use a systemd service (started at boot or via \texttt{SYSTEMD\_WANTS} when the device appears).
\end{enumerate}
